\documentclass[11pt,a4paper]{article}
\usepackage[margin=23mm,headheight=14pt,headsep=8.5mm,footskip=10mm]{geometry}
\usepackage{amsmath,amssymb,amsthm}
\usepackage[no-math]{fontspec}
\usepackage{xeCJK}
\setmainfont{Latin Modern Roman}
\setCJKmainfont{FandolSong-Regular.otf}[Script=CustomDefault,Path=assets/fonts/fandol/,BoldFont=FandolSong-Bold.otf,ItalicFont=FandolSong-Regular.otf,BoldItalicFont=FandolSong-Bold.otf]
\setCJKsansfont{FandolSong-Regular.otf}[Script=CustomDefault,Path=assets/fonts/fandol/,BoldFont=FandolSong-Bold.otf]
\xeCJKsetup{PunctStyle=quanjiao,CheckSingle=true}
\usepackage{setspace}
\setstretch{1.12}
\setlength{\parindent}{2em}
\setlength{\parskip}{3pt}
\usepackage{titlesec}
\titleformat{\section}{\centering\normalfont\Large\bfseries}{\thesection.}{0.6em}{}
\titleformat{\subsection}{\normalfont\large\bfseries}{}{0pt}{}
\titlespacing*{\section}{0pt}{1.8\baselineskip}{1.0\baselineskip}
\titlespacing*{\subsection}{0pt}{1.1\baselineskip}{0.5\baselineskip}
\usepackage{enumitem}
\setlist[enumerate]{leftmargin=2em,labelsep=0.5em,itemsep=3pt,parsep=0pt,topsep=6pt}
\providecommand{\tightlist}{\setlength{\itemsep}{3pt}\setlength{\parskip}{0pt}}
\usepackage{needspace}
\newenvironment{statementblock}{\par\Needspace{9\baselineskip}\noindent}{\par\nopagebreak[4]}
\newenvironment{proofblock}{\par\Needspace{4\baselineskip}\noindent}{\par}
\usepackage{fancyhdr}
\pagestyle{fancy}
\fancyhf{}
\fancyhead[L]{\small 数学手记}
\fancyhead[R]{\small 实变函数论 · 第二章}
\fancyfoot[C]{\thepage}
\renewcommand{\headrulewidth}{0.4pt}
\usepackage[hidelinks,unicode,bookmarksnumbered=true,linktoc=all]{hyperref}
\usepackage{bookmark}
\bookmarksetup{depth=2}
\hypersetup{pdftitle={实变函数讲义：第二章},pdfauthor={Vincent Arqain},pdfsubject={测度理论}}
\renewcommand{\contentsname}{目录}
\setcounter{tocdepth}{1}
\setcounter{secnumdepth}{1}
\widowpenalty=10000
\clubpenalty=10000
\displaywidowpenalty=10000
\predisplaypenalty=10000
\interdisplaylinepenalty=10000
\raggedbottom

\usepackage{tabularx,booktabs,array}
\renewcommand{\thesection}{2.\arabic{section}}
\emergencystretch=2em
\begin{document}
\begin{center}
{\LARGE\bfseries 实变函数讲义\par}
\vspace{8pt}{\large 第二章\quad 测度理论\par}
\vspace{6pt}{\small 依徐森林、薛春华《实变函数论》\quad 教材第 98—173 页\par}
\end{center}
\vspace{8pt}
\noindent\textbf{本章的问题}\quad 知道每个区间的长度以后，哪些集合的长度随之确定？先用覆盖定义外测度，再用分割等式决定可测性。唯一性属于生成的 $\sigma$ 环，完全性决定如何处理零集的子集。最后用平移和细尺度覆盖研究正测集与零测集的结构。
\tableofcontents
\section{环上的测度、外测度、测度的延拓}
对应教材 \S{}2.1，第 98---112 页。给半开区间指定长度以后，怎样给极限过程产生的点集指定长度？困难在于可数个环集的并可能已经离开原环。教材先在较大的集类上定义外测度，再用分割等式找回可数可加性。

\subsection*{环上的可加性与极限}\label{ux73afux4e0aux7684ux53efux52a0ux6027ux4e0eux6781ux9650}

设 \(\mathcal R\) 是 \(X\) 上的环。测度 \(\mu:\mathcal R\to[0,+\infty]\) 满足 \(\mu(\varnothing)=0\)，且对两两不交的 \(E_i\in\mathcal R\)，只要 \(E=\bigcup_i E_i\in\mathcal R\)，就有 \(\mu(E)=\sum_i\mu(E_i)\)。这个关于并集的条件使定义适用于有限区间并所成的环。

若 \(A\subset B\)，把 \(B\) 分作 \(A\) 与 \(B-A\)，立即得到单调性。若 \(E,E_i\in\mathcal R\) 且 \(E\subset\bigcup_i E_i\)，令 \(D_i=(E\cap E_i)-\bigcup_{j<i}E_j\)。这些块恰好把 \(E\) 分尽，故 \(\mu(E)=\sum_i\mu(D_i)\leq\sum_i\mu(E_i)\)。即使覆盖本身的并不在环里，这个估计仍然成立。

设下述各极限集都在定义域内。若 \(E_i\uparrow E\)，将 \(E\) 写成 \(E_1\) 与各差层 \(E_{i+1}-E_i\) 的不交并，得到 \(\mu(E_i)\uparrow\mu(E)\)。若 \(E_i\downarrow E\) 且某个 \(\mu(E_k)<\infty\)，对 \(E_k-E_i\) 用前一结论再相减，得到 \(\mu(E_i)\downarrow\mu(E)\)。在自然数的计数测度下，尾集 \(\{i,i+1,\ldots\}\) 递减到空集而测度恒为无穷，恰好说明相减需要什么条件。

在 \(\sigma\) 环上，将这两种连续性用于尾交和尾并，有 \[
\mu(\varliminf E_i)\leq\varliminf\mu(E_i),\qquad
\varlimsup\mu(E_i)\leq\mu(\varlimsup E_i).
\] 后一式要求某个尾并的测度有限。若 \(\sum_i\mu(E_i)<\infty\)，尾并测度不超过相应级数的尾和，因而 \(\mu(\varlimsup E_i)=0\)。这给出处理``发生无穷多次''的基本估计。

\subsection*{外测度如何从覆盖得到}\label{ux5916ux6d4bux5ea6ux5982ux4f55ux4eceux8986ux76d6ux5f97ux5230}

记 \(\mathcal H(\mathcal R)\) 为可被可数个 \(\mathcal R\) 中的集合覆盖的一切集合。它对子集和可数并封闭。对 \(F\in\mathcal H(\mathcal R)\) 定义 \[
\mu^*(F)=\inf\left\{\sum_{i=1}^{\infty}\mu(E_i):
  E_i\in\mathcal R,\ F\subset\bigcup_{i=1}^{\infty}E_i\right\}.
\] 当 \(F\in\mathcal R\) 时，用 \(F\) 自己覆盖得 \(\mu^*(F)\leq\mu(F)\)；前段的覆盖估计给出反向不等式。所以扩大的定义域保留了原来的数值。

外测度单调且次可数可加。证明后一条时，若 \(\sum_i\mu^*(F_i)<\infty\)，为每个 \(F_i\) 选取费用不超过 \(\mu^*(F_i)+\varepsilon 2^{-i}\) 的覆盖。把这些覆盖并成一列，再令 \(\varepsilon\downarrow0\) 即得结论；右端无穷时不等式自动成立。精度按指标分配，才能在合并覆盖以后仍只损失一个 \(\varepsilon\)。

\subsection*{分割判据与完全延拓}\label{ux5206ux5272ux5224ux636eux4e0eux5b8cux5168ux5ef6ux62d3}

称 \(E\in\mathcal H(\mathcal R)\) 为 \(\mu^*\) 可测集，如果每个 \(F\in\mathcal H(\mathcal R)\) 都满足 \[
\mu^*(F)=\mu^*(F\cap E)+\mu^*(F-E). \qquad (1)
\] 记这类集合的全体为 \(\mathcal R^*\)。次可加性已给出左端不大于右端。若 \(E\in\mathcal R\)，把 \(F\) 的任意环覆盖逐块沿 \(E\) 切开；原覆盖的总费用等于两组费用之和，因而不小于右端。取下确界即得式（1），故 \(\mathcal R\subset\mathcal R^*\)。

\textbf{延拓定理（2.1.4、2.1.5）}　\(\mathcal R^*\) 是 \(\sigma\) 环，\(\mu^*\) 在其上是完全测度，且在 \(\mathcal R\) 上等于 \(\mu\)。

\textbf{证明}　若 \(A,B\) 满足式（1），先用 \(A\) 切 \(F\)，再用 \(B\) 切所得两块。四块的外测度之和等于 \(\mu^*(F)\)。把四块按属于或不属于 \(A\cup B\) 归并，次可加性给出所需的反向不等式；对 \(A-B\) 同理。于是可测集成环。

对两两不交的可测集 \(A_i\)，令 \(S_n=\bigcup_{i\le n}A_i\)，\(S=\bigcup_i A_i\)。对前 \(n\) 块依次分割，余项由单调性估计，得到 \[
\mu^*(F)\geq\sum_{i=1}^n\mu^*(F\cap A_i)+\mu^*(F-S).
\] 令 \(n\to\infty\)，再用 \(\sum_i\mu^*(F\cap A_i)\geq\mu^*(F\cap S)\)，可知 \(S\) 可测。一般可数并先改写为不交并。最后取 \(F=S\)，上式及次可加性给出可数可加性。若 \(\mu^*(N)=0\) 且 \(A\subset N\)，则对任何 \(F\)，\(\mu^*(F\cap A)=0\) 且 \(\mu^*(F-A)\leq\mu^*(F)\)，故 \(A\) 满足式（1）。证毕。

\textbf{检验判据的反例}　取 \(\mathcal R=\{\varnothing,X\}\)，\(\mu(X)=1\)，其中 \(X\) 至少有两个点。每个非空子集的外测度均为 \(1\)。用 \(F=X\) 检验，每个非平凡子集都把 \(1\) 切成 \(1+1\)，故 \(\mathcal R^*=\mathcal R\)。用同一个 \(\mu^*\) 再作覆盖延拓也不会改变外测度：新覆盖的每块都能换成原环覆盖，并使总误差小于任意指定的 \(\varepsilon\)；反向不等式由原覆盖仍可用得到。这解释了为什么扩大可测域与在更大集类上计算外测度是两件不同的事。

\section{\texorpdfstring{$\sigma$}{σ} 有限测度、测度延拓的惟一性定理}
对应教材 \S{}2.2，第 113---122 页。上一节给出了一个延拓。现在要判断，给定环上的测度能否决定生成的 \(\sigma\) 环上每个集合的测度，以及完全化又加入了哪些集合。

\subsection*{有限性要在哪个集类上检验}\label{ux6709ux9650ux6027ux8981ux5728ux54eaux4e2aux96c6ux7c7bux4e0aux68c0ux9a8c}

教材称 \(\mu\) 在环 \(\mathcal R\) 上\textbf{有限}，是指每个 \(A\in\mathcal R\) 的测度都有限；当 \(\mathcal R\) 是代数且 \(\mu(X)<\infty\) 时称\textbf{全有限}。每个环集都能被可数个有限测度的\textbf{原环集}覆盖时，称 \(\mu\) 在 \(\mathcal R\) 上 \textbf{\(\sigma\) 有限}；若 \(\mathcal R\) 是代数，只须检查全集的这种覆盖，称全 \(\sigma\) 有限。例如有界半开区间的有限并均有有限长度，整个实轴却有无限长度。不可数集上定义在幂集上的计数测度甚至不具备 \(\sigma\) 有限性，因为可数个有限集的并至多可数。

\subsection*{为什么唯一性先在有限块上证明}\label{ux4e3aux4ec0ux4e48ux552fux4e00ux6027ux5148ux5728ux6709ux9650ux5757ux4e0aux8bc1ux660e}

\textbf{定理 2.2.1}　设 \(\mu_1,\mu_2\) 是 \(\mathcal R_\sigma(\mathcal R)\) 上的测度，在 \(\mathcal R\) 上相等，且共同的限制在 \(\mathcal R\) 上 \(\sigma\) 有限。则它们在 \(\mathcal R_\sigma(\mathcal R)\) 上相等。

\textbf{证明}　固定 \(A\in\mathcal R\) 且 \(\mu_1(A)=\mu_2(A)<\infty\)。把所有满足 \(\mu_1(E\cap A)=\mu_2(E\cap A)\) 的 \(E\in\mathcal R_\sigma(\mathcal R)\) 收在一起。该类包含原环；单调递增时用下连续性，单调递减时交集总在有限测度的 \(A\) 内，故能用上连续性。由环的单调类定理，该等式对整个生成的 \(\sigma\) 环成立。

可被可数个原环有限测度集覆盖的集合在 \(\mathcal H(\mathcal R)\) 中构成对差、可数并封闭的类，且由假设包含 \(\mathcal R\)。故每个待测的 \(E\in\mathcal R_\sigma(\mathcal R)\) 均有这种覆盖 \(A_i\)。令 \(D_i=A_i-\bigcup_{j<i}A_j\)，它们仍是有限测度的原环集。对各 \(D_i\) 使用已证等式并求和，得到 \(\mu_1(E)=\mu_2(E)\)。证毕。

\textbf{反例 2.2.1 的要点}　在 Borel 集上分别数 \(E\cap\mathbb Q\) 和 \(E\cap(\sqrt2\,\mathbb Q)\) 的点数。两个测度在非空的有界半开区间上都取无穷，在这些区间生成的环 \(\mathcal R_0\) 上相等，而对 \(\{\sqrt2\}\) 分别给出 \(0\) 与 \(1\)。它们在 Borel 域上各自都是 \(\sigma\) 有限的（可把被数的点逐一列出并加上其补集）。失去唯一性的原因正是原环里没有非空的有限测度块。

\subsection*{等测包与完成化}\label{ux7b49ux6d4bux5305ux4e0eux5b8cux6210ux5316}

记 \(\mathcal N_{\mu^*}\) 为所有外测度零集。若 \(E\in\mathcal H(\mathcal R)\) 且 \(\mu^*(E)<\infty\)，对每个 \(k\) 取一列覆盖，其并 \(G_k\in\mathcal R_\sigma(\mathcal R)\)，费用小于 \(\mu^*(E)+2^{-k}\)。令 \(G=\bigcap_k G_k\)，则 \(G\supset E\) 且 \(\mu^*(G)=\mu^*(E)\)。此时 \(G\) 称为等测包。若 \(E\) 还可测，则 \(G-E\) 的测度为零；对任意集合，仅凭等测不能推出这一步。

\textbf{结构定理（2.2.3---2.2.5）}　在原环上 \(\sigma\) 有限时，\(E\in\mathcal R^*\) 当且仅当存在 \(B,C\in\mathcal R_\sigma(\mathcal R)\) 使 \[
B\subset E\subset C,\qquad \mu^*(C-B)=0. \qquad (2)
\] 特别地，\(E\) 可写成生成的 \(\sigma\) 环中的集合添上一个外零集，也可写成其中一个集合减去一个外零集。

\textbf{证明}　用原环中的可数个有限测度块覆盖 \(E\)，在每块上对可测的交集取等测包，再把包取并，得到 \(C\supset E\) 且 \(\mu^*(C-E)=0\)。对 \(C-E\) 再取外测度为零的生成域包 \(Z\)，置 \(B=C-Z\)，即得式（2）。反过来，\(E=B\cup(E-B)\)，后项是零集的子集；由完全性，\(E\) 可测且与 \(B,C\) 等测。证毕。

若原域已是 \(\sigma\) 环，可直接把所有原域零集的子集加入。令 \(\mathcal N\) 为这些子集构成的族，写 \(E=(F\cup N)-M\)，其中 \(F\in\mathcal R\)，\(N,M\in\mathcal N\)，并定义 \(\mu'(E)=\mu(F)\)。两种表示的 \(F\) 的对称差被原域零集覆盖，所以定义与表示无关。对可数并，表示中的全部误差仍落在一个原域零集内；处理不交并时可在此零集外逐块相加。这给出最小完全延拓（定理 2.2.6）。在 \(\sigma\) 有限条件下它与 \(\mathcal R^*\) 上的延拓一致。没有该条件时两者可能不同：取不可数 \(X\) 上的可数或余可数代数及计数测度，它已完全；但每个不可数集的外测度都为无穷，任意子集均满足分割判据，故后者的定义域是整个幂集。

\section{Lebesgue 测度、Lebesgue--Stieltjes 测度}
对应教材 \S{}2.3，第 123---147 页。从区间长度出发，得到 \(\mathbb R^n\) 上的 Lebesgue 测度；将长度换成单调函数的增量，便得到可以在单点上有质量的测度。

\subsection*{从半开区间出发}\label{ux4eceux534aux5f00ux533aux95f4ux51faux53d1}

记 \(\mathcal R_0\) 为有界半开区间 \((a,b]\) 的有限并所成的环。对不交分解 \(E=\bigcup_{j=1}^k(a_j,b_j]\)，定义 \(m(E)=\sum_j(b_j-a_j)\)。两种分解的全部端点给出共同细分；长度在每一处切口相加，故定义与分解无关，且有限可加。半开约定使相邻两段在拼接处恰好只计一次。

\textbf{从有限可加到可数可加}　设 \((a,b]\) 分成可数个互不相交的 \((a_i,b_i]\)。由任意有限个小段的总长不超过 \(b-a\)，得 \(\sum_i(b_i-a_i)\leq b-a\)。反过来，把各小段右端略向右放宽，额外长度之和小于 \(\varepsilon\)。所得开区间覆盖紧集 \([a+\eta,b]\)，有有限子覆盖，故 \(b-a-\eta\leq\sum_i(b_i-a_i)+\varepsilon\)。令 \(\eta,\varepsilon\downarrow0\)。对有限区间并只须在其各分量分别使用此结论。于是 \S{}2.1 的构造适用。

由于实轴是可数个有界区间的并，\(\mathcal H(\mathcal R_0)=2^{\mathbb R}\)，且 \(\mathcal R_\sigma(\mathcal R_0)=\mathcal B\) 是 Borel 代数。用 \(m^*\) 记诱导的外测度，\(\mathcal L=\mathcal R_0^*\) 记满足分割判据的 Lebesgue 可测集类，\(m\) 记其上的完全测度。在 \(\mathbb R^n\) 中从半开长方体 \(\prod_{i=1}^n(a_i,b_i]\) 出发，指定体积 \(\prod_i(b_i-a_i)\)，同样得到 \((\mathbb R^n,\mathcal L,m)\)。平移把一个覆盖送到同费用的覆盖，反向平移给出反向不等式，所以 \(m^*(E+t)=m^*(E)\) 对任意集合成立。把分割判据整体平移，可测性也得到保留。

\subsection*{开集外逼近、闭集内逼近}\label{ux5f00ux96c6ux5916ux903cux8fd1ux95edux96c6ux5185ux903cux8fd1}

\textbf{定理 2.3.1}　集合 \(E\subset\mathbb R^n\) 可测，当且仅当对每个 \(\varepsilon>0\) 都能找到开集 \(G\supset E\) 使 \(m^*(G-E)<\varepsilon\)；等价地，可找到闭集 \(F\subset E\) 使 \(m^*(E-F)<\varepsilon\)。还等价于存在 \(F_\sigma\) 集 \(K\) 和 \(G_\delta\) 集 \(H\)，使 \[
K\subset E\subset H,\qquad m(H-K)=0. \qquad (3)
\]

\textbf{证明}　先设 \(E\) 可测且有限测度。由外测度定义取总费用接近 \(m(E)\) 的半开长方体覆盖，每个长方体稍加扩大成为开长方体，额外体积的总和可任意小。因此有开集 \(G\supset E\) 满足 \(m(G)<m(E)+\varepsilon\)；可测性使差额正好是 \(m(G-E)\)。无界或无限测度时，对 \(E\cap[-k,k]^n\) 分别给误差 \(\varepsilon 2^{-k-1}\)，最后取并。对余集作同样的开逼近再取补即得闭逼近。让两侧误差趋于零，外侧取交、内侧取并便得式（3）。若只知开逼近，对 \(H=\bigcap_k G_k\) 有 \(m^*(H-E)=0\)，于是 \(E\) 可测；只知闭逼近同理。式（3）也因完全性立即推出可测性。证毕。

注意等测包与零测差的区别。对任意 \(E\) 都可取 \(G_\delta\) 包 \(H\) 使 \(m(H)=m^*(E)\)。若 \(m^*(E)<\infty\)，可测性恰能由 \(m^*(H-E)=0\) 判别；仅有两个数相等，不能把 \(H-E\) 的外测度写成它们的差。无限测度时，即便两集可测，\((-\infty,0)\subset(-\infty,1)\) 也说明等测未必意味着相差零集。

\subsection*{Borel 集、零集与不可测集}\label{borel-ux96c6ux96f6ux96c6ux4e0eux4e0dux53efux6d4bux96c6}

由式（3），每个 Lebesgue 可测集都可由 Borel 集添去零集得到。这个扩大是严格的。三分 Cantor 集 \(C\) 有连续统势 \(\mathfrak c\)，但 \(m(C)=0\)，故其 \(2^{\mathfrak c}\) 个子集全部可测；Borel 集只有 \(\mathfrak c\) 个（第一章），所以其中有非 Borel 集。

\textbf{定理 2.3.3}　每个正测度的 Lebesgue 可测集 \(E\subset\mathbb R\) 都含有不可测子集。证明时先截取有界的正测部分 \(E_0\subset[-M,M]\)。按 \(x-y\in\mathbb Q\) 分类，在 \(E_0\) 遇到的每类中选一个代表组成 \(V\)。各有理平移 \(V+q\) 两两不交。若 \(m(V)=0\)，全部有理平移的并将以零测集覆盖 \(E_0\)。若 \(m(V)>0\)，取 \(q\in\mathbb Q\cap[0,1]\)，无穷多个平移的并具有无穷测度，却包含在 \([-M,M+1]\) 中。两种情形均矛盾。

第一章的 Cantor 函数 \(c\) 还给出了一个具体的搬运方法。\(h(x)=(x+c(x))/2\) 是 \([0,1]\) 上的递增同胚。\(c\) 在 \(C\) 的每个补区间上常值，故 \(h\) 将各补区间的长度减半，\(m(h(C))=1/2\)。在 \(h(C)\) 中选不可测子集 \(V\)，则 \(h^{-1}(V)\subset C\) 可测，却不能是 Borel 集：同胚及其逆均保持 Borel 集。这也解释了为何用连续映射讨论 Lebesgue 可测性时必须检查对零集的作用。

\subsection*{函数的跳跃怎样成为单点质量}\label{ux51fdux6570ux7684ux8df3ux8dc3ux600eux6837ux6210ux4e3aux5355ux70b9ux8d28ux91cf}

设 \(g:\mathbb R\to\mathbb R\) 单调不减且右连续。定义 \(m_g((a,b])=g(b)-g(a)\)。在原环上按不交分解求和，望远镜相消保证良定义。上一段紧覆盖论证仍可用：由右连续性选 \(\eta_i>0\) 使 \(g(b_i+\eta_i)-g(b_i)<\varepsilon 2^{-i}\)，又可使 \(g(a+\eta)-g(a)<\varepsilon\)。用有限覆盖估计增量即得可数可加性。延拓到 Carathéodory 可测域所得的完全测度称 Lebesgue--Stieltjes 测度，其在 Borel 集上的取值由这些区间增量唯一决定。

对此测度，半开方向和连续方向是一对约定。由 \((a,a+1/n]\downarrow\varnothing\) 及局部有限性，任何这样的测度都必须满足 \(g(a+1/n)\to g(a)\)。由 \((a-1/n,a]\downarrow\{a\}\) 进一步得到 \[
\begin{aligned}
m_g(\{a\})&=g(a)-g(a-0),\\
m_g((a,b))&=g(b-0)-g(a),\\
m_g([a,b])&=g(b)-g(a-0).
\end{aligned}
\] 取 \(g(x)=x\) 得到长度；取 \(g(x)=\mathbf1_{[0,\infty)}(x)\)，全部质量集中在 \(0\)；将 Cantor 函数以两端常数延拓，所得测度没有单点质量，却在零长度集 \(C\) 上集中全部质量，因为所有补区间的增量都为零。

\section{Jordan 测度、Hausdorff 测度（选读）}
对应教材选读 \S{}2.4，第 147---164 页。两种修改区间覆盖的方式给出两个不同的问题。Jordan 理论要求用有限个区间控制边界；Hausdorff 理论允许可数覆盖，并让每块的尺度趋于零。

\subsection*{零长度集与 Jordan 可测性}\label{ux96f6ux957fux5ea6ux96c6ux4e0e-jordan-ux53efux6d4bux6027}

教材称能以有限个开区间覆盖且总长任意小的集合为\textbf{零长度集}。取闭包只给有限个区间增加端点，稍向外扩大即可得到同样的估计。因此零长度等价于闭包紧且 Lebesgue 测度为零：反向从紧闭包的可数开覆盖选出有限子覆盖。\(\mathbb Q\cap[0,1]\) 是 Lebesgue 零集，其闭包却有长度 \(1\)，所以不能用总长任意小的有限覆盖盖住它。

教材记 \(\mathcal R_J=\{E\subset\mathbb R:\partial E\text{ 为零长度集}\}\)。由于并和差的边界都包含在两边界之并中，它对有限并、差封闭，也对取补封闭，是代数。单点都在其中，其可数并 \(\mathbb Q\cap[0,1]\) 却不在其中。对 \(E\in\mathcal R_J\)，内部与闭包相差零集，\(m_J(E)=m(E^\circ)=m^*(E)\)。Cantor 集本身为零长度集，每个子集的边界也在其中，故 Jordan 可测集甚至可以不是 Borel 集。用 \(\mathcal R_J\) 作可数覆盖所诱导的外测度仍是 \(m^*\)：它包含原始区间环，而每个 Jordan 可测集又可用区间覆盖逼近其测度。

\textbf{有界集的体积判据（定理 2.4.3）}　将 \(E\subset\mathbb R^n\) 放在足够大的长方体内。其特征函数 Riemann 可积，当且仅当边界能由有限个总体积任意小的开长方体覆盖。在此情形其 Riemann 积分等于 \(m(E)\)。

\textbf{证明}　对长方体作网格分割，Darboux 上下和的差就是同时遇到 \(E\) 和其补集的格子总体积。若差很小，则这些格子连同网格面覆盖边界；有限个网格面的体积为零，可把所有覆盖块略扩成开长方体。反之，给定边界的小体积开覆盖，边界紧，故它与该覆盖并集的补集有正距离（边界为空时显然）。网格直径足够小时，一切遇到边界的格子都在该并集内。每个振幅为 \(1\) 的格子必遇到边界，故上下和之差可任意小。内外格子的体积同时夹住 \(m(E)\)，得最后一项结论。证毕。

\subsection*{从覆盖的幂次读出维数}\label{ux4eceux8986ux76d6ux7684ux5e42ux6b21ux8bfbux51faux7ef4ux6570}

对 \(E\subset\mathbb R\)、\(\alpha,\delta>0\)，取开区间可数覆盖，定义 \[
\begin{aligned}
H^*_{\alpha,\delta}(E)&=\inf\left\{\sum_i |I_i|^\alpha:
  E\subset\bigcup_i I_i,\ |I_i|<\delta\right\},\\
H^*_\alpha(E)&=\lim_{\delta\downarrow0}H^*_{\alpha,\delta}(E),\qquad
d_H(E)=\inf\{\alpha>0:H^*_\alpha(E)=0\}.
\end{aligned}
\] 尺度变小时可用覆盖减少，故极限存在。按覆盖分配误差即得次可数可加性；若两个集合相距正距离，取更细的尺度，每个区间至多碰到其中一个，故并集的外测度等于二者之和。这是度量外测度。Borel 集都满足其分割判据，在可测域上得到 \(\alpha\) 次 Hausdorff 测度（定理 2.4.4、2.4.5）。

在直线上，可直接验证半直线的分割判据：给 \((-\infty,a]\) 作切分时，将原覆盖按与 \(a\) 的位置分类。所有跨过 \(a\) 的区间在左、右两侧切下的部分，分别由一个长至多 \(2\delta\) 的开区间罩住（端点可略扩）；总增费至多 \(2(2\delta)^\alpha\)。先取下确界再令 \(\delta\downarrow0\)，便得到反向不等式。半直线生成全部 Borel 集。

对 \(\beta>\alpha\)，覆盖费用满足 \(H^*_{\beta,\delta}(E)\leq\delta^{\beta-\alpha} H^*_{\alpha,\delta}(E)\)。所以临界指数以下为无穷，以上为零；定义本身不决定临界点的值。\(H^*_1=m^*\)，因为区间可细分并稍放宽，总长增加任意小。每个非退化有界区间的临界指数为 \(1\)，每个可数集的维数是 \(0\)。若 \(m^*(E)>0\)，则 \(d_H(E)=1\)；反向推断须谨慎。

\subsection*{Cantor 集的维数为何等于两个对数的比}\label{cantor-ux96c6ux7684ux7ef4ux6570ux4e3aux4f55ux7b49ux4e8eux4e24ux4e2aux5bf9ux6570ux7684ux6bd4}

从 \([0,1]\) 开始，每次在每段两端各留相对长度 \(\xi\) 的闭区间，\(0<\xi<1/2\)，极限集为 \(C_\xi\)。取 \(d=\log2/\log(1/\xi)\)，则 \(2\xi^d=1\)。第 \(k\) 层有 \(2^k\) 个长 \(\xi^k\) 的区间，总长 \((2\xi)^k\to0\)，而 \(d\) 次费用为 \(1\)。把每段略扩成开区间，可知 \(H^*_d(C_\xi)\leq1\)。这一估计只给出维数上界。

\textbf{下界的证明}　把 \([0,1)\) 中每个数的二进位数字依次当作选左、右段的指令，二进有两种展开时取以零结尾的一种。对应的嵌套区间缩到 \(C_\xi\) 中一点。该映射是有限步阶梯函数的逐点极限，故为 Borel 可测；用原区间的 Lebesgue 测度定义像测度 \(\nu\)。每个第 \(k\) 层基本段的原像是长 \(2^{-k}\) 的二进区间，故它的质量正好是 \(2^{-k}\)。

若 \(\xi^k\leq |I|<\xi^{k-1}\)，第 \(k\) 层相邻基本段的空隙至少为 \((1-2\xi)\xi^{k-1}\)。因此 \(I\) 至多遇到 \(1+(1-2\xi)^{-1}\) 个基本段（取更大的整数亦可），从而 \(\nu(I)\leq C_\xi'2^{-k}\leq C_\xi' |I|^d\)。对任一足够细的覆盖求和，\(1=\nu(C_\xi)\leq C_\xi'\sum_i|I_i|^d\)。所以 \(0<H^*_d(C_\xi)<\infty\)，且 \(d_H(C_\xi)=d\)。取 \(\xi=1/3\) 得到三分 Cantor 集的维数 \(\log2/\log3\)。

给定 \(d_k\uparrow1\) 且 \(0<d_k<1\)，取 \(\xi_k=2^{-1/d_k}\)。集合 \(\bigcup_k C_{\xi_k}\) 仍是 Lebesgue 零集，却包含维数任意接近 \(1\) 的子集，故维数为 \(1\)。这给出了与区间同维、临界测度却为零的例子。

\textbf{第 164 页的高维记号}　书末将直线上的区间换成 \(\mathbb R^n\) 的开球 \(B\)，并明确把 \(|B|\) 记作球的体积。读这里的指数时须作换算：若 \(D=\operatorname{diam}B\)、\(\omega_n\) 为单位球体积，则 \(|B|^\alpha=(\omega_n/2^n)^\alpha D^{n\alpha}\)。所以按书中体积的 \(\alpha\) 次幂计费，对应通常按直径计费的指数 \(s=n\alpha\)。例如平面正面积集合在体积计费下的临界指数是 \(1\)，通常的 Hausdorff 维数是 \(2\)。本节对直线子集的计算中 \(n=1\)，两种指数一致。

\section{测度的典型实例和应用}
对应教材 \S{}2.5，第 165---171 页。八个实例反复用到同一个区别：正测度保证局部有足够多的点，正外测度只保证覆盖不能太便宜。下列论证分别保留所需的假设。

\subsection*{平移、重叠与差集}\label{ux5e73ux79fbux91cdux53e0ux4e0eux5deeux96c6}

\textbf{例 2.5.1、2.5.2}　若 \(E\subset\mathbb R\) 可测且 \(m(E)>0\)，它既含差为无理数的两点，也含差为非零有理数的两点。前者只要固定 \(x_0\in E\)，去掉可数集 \(x_0+\mathbb Q\)。后者先取有界正测部分 \(A\)，用 \([0,1]\) 内不同有理数 \(q_j\) 作平移。若 \(A+q_j\) 两两不交，包含它们的同一有界区间就有无限测度。故有 \(x+q_i=y+q_j\)，给出非零有理差。

\textbf{小平移引理}　若 \(A\in\mathcal L\) 且 \(m(A)<\infty\)，则 \(m((A+h)\triangle A)\to0\)（\(h\to0\)）。给定 \(\varepsilon>0\)，先选紧集 \(K\subset A\) 和开集 \(G\supset A\) 使 \(m(G-K)<\varepsilon\)。以有限个包含于 \(G\) 的有界开长方体盖住 \(K\)，其并记为 \(P\)，则 \(m(A\triangle P)<\varepsilon\)。每个长方体 \(I\) 满足 \(m((I+h)\triangle I)\to0\)，因为交集的各边长度趋于原边长。有限并也满足，故 \[
m((A+h)\triangle A)\leq 2m(A\triangle P)+m((P+h)\triangle P).
\] 先令 \(h\to0\) 再令 \(\varepsilon\downarrow0\) 即得。在直线以及 \(\mathbb R^n\) 上证明完全相同。

因此若 \(0<m(A)<\infty\)，则 \(m(A\cap(A+h))\to m(A)>0\)。对任意正测集截取这样的 \(A\) 即知，其差集包含原点的一个邻域（Steinhaus 定理）。例 2.5.4 的命题随即成立：若对每个正整数 \(n\) 均有 \(m(E\cap(E+1/n))=0\)，则 \(m(E)=0\)。这里得到的是正测重叠，比找到一对相差 \(1/n\) 的点更强。

\subsection*{外测度能分出多大的一块}\label{ux5916ux6d4bux5ea6ux80fdux5206ux51faux591aux5927ux7684ux4e00ux5757}

\textbf{例 2.5.3}　选 \S{}2.3 中有界正测集的有理差代表集 \(V\)。它有正外测度，否则全部有理平移的并仍是零集。取 \(q_j\in[0,1]\cap\mathbb Q\) 互异，各 \(V+q_j\) 两两不交且同处一个有界区间内。因此它们的外测度之和为无穷，并集的外测度却有限。这说明 \S{}2.1 的可测条件在何处发挥作用。

\textbf{例 2.5.5}　对任意 \(E\subset\mathbb R\)，若 \(0\leq q<m^*(E)\)，有 \(E_1\subset E\) 满足 \(m^*(E_1)=q\)。令 \(u(r)=m^*(E\cap(-r,r))\)，则 \(u(0)=0\)，且 \(0\leq u(s)-u(r)\leq2(s-r)\)（\(s>r\)），后者由两个边缘区间覆盖给出。以有界可测环带切开 \(E\)，分割判据给出前 \(k\) 带的外测度等于各带外测度之和；取极限并用次可数可加性，得到 \(u(r)\uparrow m^*(E)\)。由介值定理取 \(u(r)=q\) 即可。无需假定被切的 \(E\) 可测。

\textbf{例 2.5.6}　对抽象外测度 \(\mu^*\)（由环上测度诱导），有 \[
\mu^*(A\cup B)+\mu^*(A\cap B)\leq\mu^*(A)+\mu^*(B).
\] 右端有限时，分别给 \(A,B\) 取生成 \(\sigma\) 环中的等测包，对这两个可测包用容斥，再由单调性缩回；右端无穷时显然。若 \(A\) 可测，分别沿 \(A\) 切开 \(A\cup B\) 与 \(B\) 就得到等式，全程无需作无穷量相减。严格不等式可取不可测 \(A\subset[0,1]\) 及 \(B=[0,1]-A\)：如有 \(m^*(A)+m^*(B)=1\)，取两者的 Borel 等测包，再与 \([0,1]\) 相交，则两包之交为零集，\(A\) 夹在相差零集的两个 Borel 集之间，矛盾。

\subsection*{稠密平移不变与每段都两边有量}\label{ux7a20ux5bc6ux5e73ux79fbux4e0dux53d8ux4e0eux6bcfux6bb5ux90fdux4e24ux8fb9ux6709ux91cf}

\textbf{例 2.5.7}　设 \(D\) 在直线上稠密，\(E\) 可测，且每个 \(h\in D\) 都满足 \(m((E+h)\triangle E)=0\)。则 \(m(E)=0\) 或 \(m(E^c)=0\)。

\textbf{证明}　固定 \(\ell>0\)，令 \(v(t)=m(E\cap(t,t+\ell))\)。两个区间的对称差长度不超过 \(2|s-t|\)，故 \(v\) 连续。每个 \(h\in D\) 都是 \(v\) 的周期，由稠密性和连续性，\(v\) 为常数。于是区间内 \(E\) 的测度只依赖于区间长度。把区间等分并用长度的连续性，得对所有有界区间 \(I\) 有 \(m(E\cap I)=c\,m(I)\)，其中 \(c=m(E\cap(0,1))\)。两边作为 Borel 测度由 \S{}2.2 的唯一性在 Borel 集上一致；任意 Lebesgue 可测集与某个 Borel 集相差零集，两边仍一致。取 \(A=E\cap[-k,k]\) 得 \((1-c)m(A)=0\)。若某个 \(m(A)>0\)，则 \(c=1\)，每个有界区间中 \(E^c\) 都为零测，故 \(m(E^c)=0\)。证毕。

\textbf{例 2.5.8}　要使每个非退化子区间内 \(E\) 与其补集都占正测度，不能在任何一边留下整段开区间。列出 \((0,1)\) 内所有有理端点开区间 \(I_j\)。在 \(I_1\) 中放两个不交的正测度疏朗紧集 \(K_1,L_1\)。第 \(j\) 次，先删掉此前有限多个这样的紧集，\(I_j\) 仍含开区间，在其中再放 \(K_j,L_j\)。所需紧集可在任一闭小区间中挖去覆盖其有理点的开区间，总长取小于原长的一半。令 \(E=\bigcup_j K_j\)。它是 Borel 集，每个非退化区间都包含某个 \(I_j\)，其中 \(K_j\subset E\) 且 \(L_j\subset E^c\)，故两侧都有正测度。这个构造也为复习题中``正测度却没有内部''的问题提供了例子。

\clearpage\phantomsection\addcontentsline{toc}{section}{第二章习题与复习题}\section*{第二章习题与复习题}
教材第二章的三组节后练习分别有 5、10、21 题，共 36 道大题。\S{}2.4 没有单列练习，\S{}2.5 的八个例子见分节讲义。章末复习题另有 18 道，以下按原号给出问题、方法和论证。页码为印刷页；430 页扫描版的 PDF 页码须加 6。

\noindent\textbf{复习题题号}\quad \hyperlink{review-q1}{1}\quad  \hyperlink{review-q2}{2}\quad  \hyperlink{review-q3}{3}\quad  \hyperlink{review-q4}{4}\quad  \hyperlink{review-q5}{5}\quad  \hyperlink{review-q6}{6}\quad  \hyperlink{review-q7}{7}\quad  \hyperlink{review-q8}{8}\quad  \hyperlink{review-q9}{9}\quad  \hyperlink{review-q10}{10}\quad  \hyperlink{review-q11}{11}\quad  \hyperlink{review-q12}{12}\quad  \hyperlink{review-q13}{13}\quad  \hyperlink{review-q14}{14}\quad  \hyperlink{review-q15}{15}\quad  \hyperlink{review-q16}{16}\quad  \hyperlink{review-q17}{17}\quad  \hyperlink{review-q18}{18}

\subsection*{练习 2.1 · 测度与外测度（第 112 页）}\label{ux7ec3ux4e60-2.1-ux6d4bux5ea6ux4e0eux5916ux6d4bux5ea6ux7b2c-112-ux9875}

\textbf{1}　对上极限、极限和尾和中的有限性条件，分别考察计数测度的递减尾集，以及各项测度趋于零而尾和仍发散的区间列。例如在 \([0,1)\) 逐层列出二进分割的全部区间：单项长度趋于零，每点仍出现无穷多次。

\textbf{2}　以 \(2^{-n}\) 为权将测度相加。对不交集列所得双重非负级数，两种求和次序均等于有限矩形部分和的上确界。若各测度均不超过 \(1\)，权重之和即给出界。

\textbf{3}　逐集取测度的极限保留有限可加性。反例取 \(\mathbb N\) 的有限或余有限代数，在第 \(n\) 点放一单位质量；极限对有限集为零，对余有限集为一。所有单点零测，全集却不为零。

\textbf{4}　教材此处的``原子''指有正测度的单点。按 \(\mu(\{x\})\geq 1/k\) 分层，若所有环集的测度至多为 \(1\)，每层最多有 \(k\) 个点。

\textbf{5}　在 \(\sigma\) 环上，把近似最优覆盖分别取并，再对这些并取交，得到等测包；费用为无穷时任一生成域覆盖包已经等测。仅有环时可取半开区间环及单点集；该点外测度为零，却没有含它的零长度环集。

\subsection*{练习 2.2 · 延拓与内测度（第 121---122 页）}\label{ux7ec3ux4e60-2.2-ux5ef6ux62d3ux4e0eux5185ux6d4bux5ea6ux7b2c-121122-ux9875}

\textbf{1}　可数全集上，有限子集环的计数测度延拓给所有子集通常的点数。权重 \(\alpha=0\) 时所有子集均取零；\(\alpha>0\) 时按点数乘权，包括无穷点数。

\textbf{2}　令 \(X\) 至少有两点，在平凡代数上给 \(X\) 质量 \(1\)。Carathéodory 域只有平凡两集，任一指定点的单位点质量却把它延拓到整个幂集。

\textbf{3}　比较先限制环再延拓与先延拓再取迹。取实轴的半开区间环，\(E=\{0\}\)。包含在 \(E\) 中的原环集只有空集，不能覆盖 \(E\)；而原延拓在 \(E\) 上的迹包含 \(E\)。

\textbf{4}　原环上的有限块可数覆盖性质对差和可数并封闭，因而传到生成域。反向用 \S{}2.2 中数有理点的反例，比较同一个测度在两个定义域中可用的覆盖块。

\textbf{5}　原环中零测集对差与有限并封闭；若原环是 \(\sigma\) 环，则亦对可数并封闭。取有限集环上的零测度，零测集族未必为 \(\sigma\) 环；换成计数测度后该族只含空集。

\textbf{6}　在 \(\sigma\) 环上，任一环覆盖可取其并，所以外测度也可写成所有环集大包的测度下确界。反向将单个大包看成一组覆盖。

\textbf{7}　内测度 \(\mu_*(E)=\sup\{\mu(F):F\in\mathcal R,\ F\subset E\}\) 从里逼近。对不交集各取有限个内逼近块再取并，得到超可加不等式；在 \(\sigma\) 环中可先取可数并。用原环集分割时，把每个逼近块也作同一分割；另外从两边分别选块合并，得到反向不等式。

\textbf{8}　在总测度有限的代数上考察 \(\mu(X)-\mu^*(X-E)\)。对余集用外测度不等式，与第 7 题逐项比较。严格差别可取 \([0,1]\) 上的有限区间并代数（端点可单列）、长度测度，以及形如正测度疏朗紧集的 \(E\)。无非退化区间能从内逼近它，但其余集的外测度等于余集的长度。

\textbf{9}　对有限测度的 \(\sigma\) 代数，分别取趋近内、外测度的内外包，用可数并、交达到两端。数值相等时，内外包之差为零。无限总测度时，半直线与有界不可测集的不交并可使两值都为无穷而集合不可测。

\textbf{10（原题条件辨析）}　外测度等式总成立，把测试集沿不交的 \(A,B\in\mathcal R^*\) 依次切开即可。内测度等式在原题的一般条件下不成立。取 \(X=\{a,b\}\)，\(\mathcal R=\{\varnothing,X\}\)，\(\mu(X)=+\infty\)。任一非空集的外测度为无穷，两个单点都满足分割判据；但两单点的内测度均为零，取 \(E=X\)，它们的并的内测度为无穷。若增加原测度 \(\sigma\) 有限，则 \S{}2.2 的夹逼定理使每个完成域可测块都有原域中的等测内核。用 \(A,B\) 切开 \(E\cap(A\cup B)\) 中任一原域内逼近块，再取这两个内核，即得所缺的一向不等式；另一向由分别内逼近后取并得到。

\subsection*{练习 2.3 · 具体测度（第 144---147 页）}\label{ux7ec3ux4e60-2.3-ux5177ux4f53ux6d4bux5ea6ux7b2c-144147-ux9875}

\textbf{1}　欲从可测 \(E\) 中取指定测度 \(a<m(E)\) 的无内点紧集，先取测度大于 \(a\) 的紧子集，再挖去覆盖有理点的小开集，使余量仍大于 \(a\)。用半直线水平截取；累计长度随端点连续。

\textbf{2}　用分割判据处理``可测小集与大集等外测度''和容斥式。一般集合的外包取 \(G_\delta\)，内逼近取 \(F_\sigma\)；需分别达到外测度、内测度，不能据此断言内外包之差为零。

\textbf{3}　对递增集列分别取等外测包，交叠成嵌套的包，可推出外测度下连续。对递减集列用内逼近和有限性。有意交换内外测度时，可以把一个区间分成可数个分别不含正测可测子集的集合，考察有限部分并与余下的尾集；构造时还须验证有限并的内测度仍为零，例如用有理差代表集的有限个有理平移（再避开有限多个平移之差，作不交的小有理平移）。

\textbf{4}　在有界区间中删去覆盖有理点的小开集，留下只含无理点的正测闭集。题给的稠密开集 \(G\) 总长有限，若与闭集 \(F\) 几乎相等，\(G-F\) 是零测开集，必空；稠密性将迫使 \(F=\mathbb R\)，产生矛盾。

\textbf{5}　同时从外侧用开集包住 \(E\) 与 \(E^c\)，其重叠就是误差。用紧内逼近判别可测性；从可数拓扑基删去所有与 \(E\) 交成零集的球，就能找到每个邻域中都有正测度的点。对任意 \(A\) 与递减可测集的交，先给 \(A\) 取局部等测包，再用有限测度的上连续性。

\textbf{6}　在单位正方形中，\(|\sin x|<1/2\) 留下宽 \(\pi/6\) 的竖条。\(\cos(x+y)\) 取有理值的点落在可数条直线上，删去它们不改面积，故答案为 \(\pi/6\)。

\textbf{7}　下极限用递增尾交，上极限用递减尾并；后者按题设从一个有限测度的尾并开始。

\textbf{8}　有限外测度时，等测包与原集之间不能再放进正测可测块。将等测包沿任一可测集切开，所得两边同时保持等测；无限情形可分有界块构造局部等测包。对于两个集合并可测且外测度之和等于并的有限测度，先取各自的可测包限制在并内；两包的交为零，两集合即分别可测。

\textbf{9}　满测度集取可数交可对补集求并；正下界由余集测度相加得到。区间覆盖小问可先作有限并，逐次剔除被其余区间覆盖的区间；不冗余的有限族覆盖重数至多二。对并集的补误差计数，再令有限部分并递增。

\textbf{10}　若零集 \(Z\) 的像为正测度，就从 \(f(Z)\) 中取不可测子集 \(V\)。集合 \(Z\cap f^{-1}(V)\) 是零集的子集，按假设其像 \(V\) 应可测，矛盾。

\textbf{11}　将有界区间内的覆盖重复细化。每轮的总长度至多乘以同一个 \(q<1\)，极限给出零外测度。

\textbf{12}　四小问分别考察每个区间固定占半、点差碰撞、十进数码限制、不可测集与大测度集之交。第一问用正则逼近得到占比可任意接近一的区间；第二问比较 \(E+x_i\) 在长 \(2\) 的区间内能否互不重叠（所给点须互异才保证距离非零）。第三问固定缺失的数码，每层至多留下 \(9^k\) 个长 \(10^{-k}\) 的区间，再处理有限种缺失数码及双重展开；势的下界可用只含给定数码的无限尾数编码。第四问比较不可测集的内、外测度缺口与删去部分的测度。

\textbf{13}　\(E\) 与关于 \(a\) 的反射像合起来覆盖 \((a-\delta,a+\delta)\)，反射保持测度，故 \(m(E)\geq\delta\)。

\textbf{14}　用 \(W=(W\triangle E)\triangle E\)，检验对称差与可测域的封闭性。

\textbf{15}　把任意测试集和它被 \(T(E)\) 分出的两块拉回，以外测度不变性使用分割判据。若``一一''仅指单射，可先在像域中操作；像域可测性需要另外证明。按一一对应理解时可直接用逆映射。

\textbf{16}　\(A_k\subset E\subset B_k\)，且 \(m(B_k-A_k)\to0\)。用 \(\bigcup_k A_k\) 与 \(\bigcap_k B_k\) 夹住 \(E\)，两者之差外测度为零。

\textbf{17}　取 \(\mathbb Q\cap[0,1]\)。任一包含它的闭集含整个 \([0,1]\)，故一般的开外逼近不能换为闭外逼近。

\textbf{18}　给 \([0,1]\) 内所有有理点取很短的开邻域，构造在 \((0,1)\) 中稠密且长度小于 \(1/2\) 的开集；闭包的测度为 \(1\)。

\textbf{19}　若补集有正测度，可在某段上使补集占比任意接近一，这与原集在所有子区间的统一正比例下界矛盾。

\textbf{20}　任意点集都是开集族的交：把每个不在其中的单点删去。因此任意交甚至可能不可测。可数交才由 Borel 代数保证封闭。

\textbf{21}　要从矩形推广旋转不变性，可先在开集上作可数个近乎不交的方形分解。旋转后的方形仍有原面积（Jordan 体积），边界零测；由开包逼近得到任意集合的外测度在旋转下不变。可测性再由分割判据推出。

\subsection*{复习题 2 · 覆盖、平移与零集（第 171---172 页）}\label{ux590dux4e60ux9898-2-ux8986ux76d6ux5e73ux79fbux4e0eux96f6ux96c6ux7b2c-171172-ux9875}

\subsubsection*{问题 1　子列的公共部分}\hypertarget{review-q1}{}\label{ux95eeux9898-1-ux5b50ux5217ux7684ux516cux5171ux90e8ux5206}

设 \(E_n\subset[0,1]\) 可测。若 \(m(E_n)\to1\)，给定 \(0<\alpha<1\)，依次取 \(n_k>n_{k-1}\) 使 \(m(E_{n_k}^c)<(1-\alpha)2^{-k-1}\)（补集在 \([0,1]\) 中取）。由次可加性，交集测度至少为 \(1-(1-\alpha)/2>\alpha\)。若只有某子列的测度趋于 \(1\)，在该子列中作同样选择即可。

统一正下界则不足。令 \(E_n\) 为 \([0,1)\) 中二进展开第 \(n\) 位为 \(1\) 的点集；双重展开取以零结尾的表示。指定 \(k\) 个位置的数字时，将区间分到最大位数对应的层，每种指定模式占总长 \(2^{-k}\)。所以 \(m(E_n)=1/2\)，任意无限子列的交集测度却为零。这些 \(E_n\) 彼此不同。

\subsubsection*{问题 2　不可测集与其补集的外包}\hypertarget{review-q2}{}\label{ux95eeux9898-2-ux4e0dux53efux6d4bux96c6ux4e0eux5176ux8865ux96c6ux7684ux5916ux5305}

设不可测 \(E\subset\mathbb R^n\)。如果不存在题设的统一正下界，就有可测 \(A_k\supset E\)、\(B_k\supset E^c\) 满足 \(m(A_k\cap B_k)<2^{-k}\)。令 \(A=\bigcap_k A_k\)、\(B=\bigcap_k B_k\)，则 \(m(A\cap B)=0\)。由于 \(B^c\subset E\subset A\) 且两侧相差 \(A\cap B\)，完全性将推出 \(E\) 可测，矛盾。

\subsubsection*{问题 3　任意角集的扇面外测度}\hypertarget{review-q3}{}\label{ux95eeux9898-3-ux4efbux610fux89d2ux96c6ux7684ux6247ux9762ux5916ux6d4bux5ea6}

设 \(E\subset(-\pi,\pi)\)，扇面由 \(a<r<b\) 和 \(\theta\in E\) 确定。用开区间覆盖角集，先与 \((-\pi,\pi)\) 相交。每段长 \(\ell\) 的角区间生成一个面积为 \((b^2-a^2)\ell/2\) 的普通扇环；边界圆弧和线段零测。对覆盖的各段求和、再取下确界即得 \(m^*(S)\leq (b^2-a^2)m^*(E)/2\)。角集无需可测。

\subsubsection*{问题 4　高占比方体与差集}\hypertarget{review-q4}{}\label{ux95eeux9898-4-ux9ad8ux5360ux6bd4ux65b9ux4f53ux4e0eux5deeux96c6}

设 \(m(E)>0\)。截取有限正测部分 \(A\)，选开集 \(G\supset A\) 使 \(\lambda m(G)<m(A)\)，其中 \(0<\lambda<1\)。用半开二进方体分解 \(G\)：从边长不超过一且闭包在 \(G\) 内的方体中取极大者。每点都在某个这样的方体内，二进方体或嵌套或不交，故这些极大方体可数、不交且盖满 \(G\)。若每块 \(Q\) 均有 \(m(E\cap Q)\leq\lambda m(Q)\)，求和就与 \(m(A)>\lambda m(G)\) 矛盾。将找到的方体改为开方体只改变零测边界。

差集包含零点邻域可直接应用 \S{}2.5 的小平移引理。也可取前问中 \(m(E\cap Q)>3m(Q)/4\) 的方体，令 \(A=E\cap Q\)。对足够小的 \(h\)，\(m(Q\cup(Q+h))<3m(Q)/2\)。若 \(A\) 与 \(A+h\) 不交，它们之并的测度将大于这个界，故 \(h\in A-A\subset E-E\)。

\subsubsection*{问题 5　和集与中点封闭}\hypertarget{review-q5}{}\label{ux95eeux9898-5-ux548cux96c6ux4e0eux4e2dux70b9ux5c01ux95ed}

正测可测集 \(E\) 有 \(E+E\) 包含一个开区间，由第 16 题对 \(A=B=E\) 的证明即得。若还对中点封闭，则 \((E+E)/2\subset E\)，将该开区间缩小一半便落在 \(E\) 中。差集定理本身给出的是 \(E-E\)，若改用它，须先作反射并证明适当的交集为正测。

\subsubsection*{问题 6　过快的有理逼近}\hypertarget{review-q6}{}\label{ux95eeux9898-6-ux8fc7ux5febux7684ux6709ux7406ux903cux8fd1}

固定 \(M\in\mathbb N\)，只看 \(x\in[-M,M]\)。对分母 \(q\)，能出现的分子有至多 \(2(M+1)q+1\) 个，每个贡献一个长 \(2q^{-\alpha}\) 的区间。因 \(\alpha>2\)，这些覆盖在 \(q\geq N\) 处的总长度不超过常数乘以 \(\sum_{q\geq N}q^{1-\alpha}\to0\)。有无穷多个互素分数逼近的点必在每个尾覆盖中，故它们在此盒内为零集；再令 \(M\) 遍历自然数。上界估计中允许非互素分数，只会扩大覆盖。

\subsubsection*{问题 7　有理差代表集}\hypertarget{review-q7}{}\label{ux95eeux9898-7-ux6709ux7406ux5deeux4ee3ux8868ux96c6}

令 \(W\) 在 \(\mathbb R^n/\mathbb Q^n\) 的每个等价类中恰取一点。若 \(W\) 可测且正测，由第 4 题，\(W-W\) 含零点邻域，故含非零有理向量，违背代表的唯一性。若 \(W\) 可测且零测，可数个平移 \(W+q\) 的并仍零测，却覆盖 \(\mathbb R^n\)。因此 \(W\) 不可测。

\subsubsection*{问题 8　单调光滑映射的原像}\hypertarget{review-q8}{}\label{ux95eeux9898-8-ux5355ux8c03ux5149ux6ed1ux6620ux5c04ux7684ux539fux50cf}

设 \(f\in C^1(\mathbb R)\) 且 \(f'>0\)。对 \([-k,k]\) 有 \(c_k=\min f'>0\)，中值定理给出该区间像上的逆函数为 \(1/c_k\)-Lipschitz。Lipschitz 函数把零集映成零集，因为可将区间覆盖的总长至多放大 Lipschitz 常数倍。若 \(E\) 可测，将它写成 Borel 集与零集之并。前者由连续性有 Borel 原像；后者在每个 \([-k,k]\) 内的原像由上述逆映射送出，故为零集。可数取并，得到 \(f^{-1}(E)\) 可测。

\subsubsection*{问题 9　稠密平移决定 Borel 测度}\hypertarget{review-q9}{}\label{ux95eeux9898-9-ux7a20ux5bc6ux5e73ux79fbux51b3ux5b9a-borel-ux6d4bux5ea6}

设 \(\mu\) 对紧集有限，题设半开区间的质量在稠密集 \(D\) 给出的平移下不变。先消除端点：对任意 \(x\)，令区间 \((x-1/k,x]\) 递减到 \(\{x\}\)，由上连续性知 \(\mu(\{x+h\})=\mu(\{x\})\) 对 \(h\in D\) 成立。若有一个正质量点，选无穷多个有界的不同 \(h\in D\)，会在某个紧区间内累积无穷质量，矛盾。因此无原子。

固定 \((a,b]\)。当 \(h_j\to h\)，两个平移区间的对称差只落在端点附近；由于单点零测且紧区间测度有限，这些小邻域的质量趋于零。故稠密平移不变性延伸到所有实数平移。令 \(\lambda=\mu((0,1])\)，等分单位区间先得有理长度区间的质量为长度乘 \(\lambda\)，再由下连续性推广到任意长度。在半开区间环上两测度均 \(\sigma\) 有限，由定理 2.2.1 得 \(\mu=\lambda m\) 在所有 Borel 集上成立。其中 \(\lambda=0\) 时指零测度。

\subsubsection*{问题 10　正测集上有界的加法函数}\hypertarget{review-q10}{}\label{ux95eeux9898-10-ux6b63ux6d4bux96c6ux4e0aux6709ux754cux7684ux52a0ux6cd5ux51fdux6570}

设 \(f(x+y)=f(x)+f(y)\)，且在正测集 \(E\) 上 \(|f|\leq M\)。由第 4 题，\(E-E\) 含 \((-\delta,\delta)\)，故其上 \(|f|\leq2M\)。若 \(|x|<\delta/n\)，则 \(|f(x)|=|f(nx)|/n\leq2M/n\)，从而 \(f\) 在零点连续。加法性给出各处连续和 \(f(q)=qf(1)\)（\(q\in\mathbb Q\)），取有理数逼近任意实数便得 \(f(x)=xf(1)\)。

\subsubsection*{问题 11　完全集中的零测完全集}\hypertarget{review-q11}{}\label{ux95eeux9898-11-ux5b8cux5168ux96c6ux4e2dux7684ux96f6ux6d4bux5b8cux5168ux96c6}

在与非空完全集 \(E\) 的点相交于内部的闭区间中，取两个不同的 \(E\) 点，在它们周围取不交的更小闭区间，均含于父区间内部。每个点都是聚点，故每支均可继续。第 \(k\) 层的 \(2^k\) 个区间各长至多 \(4^{-k}\)。各层并的交 \(K\) 是非空紧集，所有分支的极限点因 \(E\) 闭而属于 \(E\)；第 \(k\) 层总长至多 \(2^{-k}\)，故 \(m(K)=0\)。每个分支在任一深层都有另一条分支与它分开，故 \(K\) 无孤立点。

\subsection*{复习题 2 · 导数、像集与构造（第 172---173 页）}\label{ux590dux4e60ux9898-2-ux5bfcux6570ux50cfux96c6ux4e0eux6784ux9020ux7b2c-172173-ux9875}

\subsubsection*{问题 12　只在给定点集上控制导数}\hypertarget{review-q12}{}\label{ux95eeux9898-12-ux53eaux5728ux7ed9ux5b9aux70b9ux96c6ux4e0aux63a7ux5236ux5bfcux6570}

设 \(f\) 在 \(E\subset[a,b]\) 的各点可导，\(|f'|<M\)。固定 \(\varepsilon>0\)，令 \(E_n\) 收集满足题设一致半径条件的点：当 \(y\in[a,b]\) 且 \(|y-x|\leq1/n\) 时，\(|f(y)-f(x)|\leq(M+\varepsilon)|y-x|\)。则 \(E_n\uparrow E\)。以长小于 \(1/n\)、总长小于 \(m^*(E_n)+\varepsilon\) 的区间 \(I_j\) 覆盖 \(E_n\)（较长区间可细分并略放宽）。在 \(E_n\cap I_j\) 上任意两点的像差至多为 \((M+\varepsilon)|I_j|\)，故像的外测度不超过此数。求和得 \[
m^*(f(E_n))\leq(M+\varepsilon)(m^*(E_n)+\varepsilon).
\]

此处像集未必已知可测。所需的外测度下连续性可这样证明：若 \(S_n\uparrow S\)，分别取 Borel 等测包 \(H_j\supset S_j\)，令 \(G_n=\bigcap_{j\geq n} H_j\)，则 \(S_n\subset G_n\subset H_n\) 且 \(G_n\uparrow\)。对 \(G_n\) 用测度下连续性，得 \(m^*(S)\leq\lim m(G_n)=\lim m^*(S_n)\)，另一向由单调性。应用于 \(f(E_n)\)，再令 \(\varepsilon\downarrow0\)，结论为 \(m^*(f(E))\leq M m^*(E)\)。

\subsubsection*{问题 13　可导函数保持零集的像零测}\hypertarget{review-q13}{}\label{ux95eeux9898-13-ux53efux5bfcux51fdux6570ux4fddux6301ux96f6ux96c6ux7684ux50cfux96f6ux6d4b}

对 \(C^1\) 函数，在每个紧区间上用导数上界和中值定理，便得局部 Lipschitz 性及结论。若只知处处可导，令 \(Z_{k,j}=Z\cap[-j,j]\cap\{x:|f'(x)|<k\}\)，每个 \(Z_{k,j}\) 仍有外测度零。第 12 题给出 \(m^*(f(Z_{k,j}))=0\)，取可数并即得 \(m(f(Z))=0\)。这里无需给导数找一个邻域内的上界。

\subsubsection*{问题 14　像零测与导数几乎处处为零}\hypertarget{review-q14}{}\label{ux95eeux9898-14-ux50cfux96f6ux6d4bux4e0eux5bfcux6570ux51e0ux4e4eux5904ux5904ux4e3aux96f6}

若 \(f'=0\) 在 \(E\) 上除一零集外成立，例外集的像由第 13 题的分层证明为零（该证明只用到这些点处可导）。在余下各点对第 12 题取任意 \(M>0\)，令 \(M\downarrow0\)，其像亦零测。

反之，设 \(m(f(E))=0\)。把 \(E\) 中 \(f'(x)\ne0\) 的点按 \(k,n\) 分组，使 \(|f(y)-f(x)|\geq k^{-1}|y-x|\) 对 \(|y-x|<1/n\) 的所有定义域内 \(y\) 成立。导数非零保证这样的组覆盖全部待证点。在每一组中再按长小于 \(1/n\) 的区间切块；对块内任意两点上述估计同时有效，故 \(f\) 在该块上单射且逆为 \(k\)-Lipschitz。块的像是零集的子集，逆映射也把它送成零集（用覆盖和像的直径估计即可）。对所有块取并，得到 \(f'=0\) 在 \(E\) 上几乎处处成立。

\subsubsection*{问题 15　一维 Sard 结论}\hypertarget{review-q15}{}\label{ux95eeux9898-15-ux4e00ux7ef4-sard-ux7ed3ux8bba}

取临界点集 \(E=\{x:f'(x)=0\}\)。第 12 题对每个 \(M>0\) 都给出 \(m^*(f(E))\leq M(b-a)\)，故临界值集零测。这里只要求临界点处的导数存在并为零。

\subsubsection*{问题 16　两个正测集的和含有开区间}\hypertarget{review-q16}{}\label{ux95eeux9898-16-ux4e24ux4e2aux6b63ux6d4bux96c6ux7684ux548cux542bux6709ux5f00ux533aux95f4}

分别在 \(A\)、\(B\) 中找到占比大于 \(3/4\) 的二进区间 \(I,J\)（第 4 题的证明）。可使两者等长：将较长的二进区间细分，至少一个子段保留原占比。设共同长度为 \(\ell\)，两段中点之和为 \(s\)。则 \(I\) 与 \(s-J\) 除端点外重合。对 \(|t-s|<\ell/4\)，\(I\cup(t-J)\) 的长度小于 \(5\ell/4\)；而 \(A\cap I\) 与 \(t-(B\cap J)\) 的长度之和大于 \(3\ell/2\)，必相交。交点给出 \(t=a+b\)。整个 \((s-\ell/4,s+\ell/4)\) 都在 \(A+B\) 中。

\subsubsection*{问题 17　无正测乘积子集的平面正测集}\hypertarget{review-q17}{}\label{ux95eeux9898-17-ux65e0ux6b63ux6d4bux4e58ux79efux5b50ux96c6ux7684ux5e73ux9762ux6b63ux6d4bux96c6}

从单位正方形删去所有 \(x+y\in\mathbb Q\) 的点，余集 \(S\) 的面积为 \(1\)，因为只删去了可数条直线的一部分。若 \(A\times B\subset S\) 且 \(A,B\) 都是正测集，则 \(A+B\) 不含有理数。第 16 题却说明 \(A+B\) 含有一个区间，矛盾。

\subsubsection*{问题 18　导函数在完全集上的间断性}\hypertarget{review-q18}{}\label{ux95eeux9898-18-ux5bfcux51fdux6570ux5728ux5b8cux5168ux96c6ux4e0aux7684ux95f4ux65adux6027}

\textbf{原题的限制。} 第（1）问要求在给定的非空完全集 \(C\) 的每一点使导函数不连续。取 \(C=[0,1]\) 即见原断言不能成立：导函数在内部是连续差商的逐点极限，第一章的 Baire 定理保证它在任一内区间有连续点。对一般闭集，同一理由说明无内点是必要条件。以下对无内点的非空完全集给出构造，连续与间断均相对于函数的整个定义域而言。

\textbf{构造}　列出 \([0,1]\) 中所有与 \(C\) 相交的有理端点相对开基区间 \(U_j\)（包括端点处的相对邻域）。在每个 \(U_j-C\) 内取闭小区间 \(J_j\)，使它们两两不交，且 \(J_j\) 上每一点到 \(C\) 的距离小于 \(1/j\)。确可做到：在 \(U_j\cap C\) 取一点，它同先前有限个闭小区间有正距离，而其任意小邻域都含有补集的开小段。

取一个光滑函数 \(\phi\)，支集在 \((-1,1)\) 内，满足 \(|\phi|\leq1\)，且某点的 \(|\phi'|=1\)。在 \(J_j\) 内以中心 \(c_j\) 和半径 \(r_j\) 放入小峰 \[
u_j(x)=r_j\,\phi((x-c_j)/r_j),\qquad
0<r_j<\min\{|J_j|/2,\operatorname{dist}(J_j,C)^2\}.
\] 令 \(f=\sum_j u_j\)。每点至多遇到一峰，在 \(C\) 外局部只有有限项（各区间最终都靠近 \(C\)），故在补集上光滑。在 \(C\) 上 \(f=0\)，且 \(|f(y)|\leq\operatorname{dist}(y,C)^2\)，所以在那里导数为零。每个 \(C\) 点的任意邻域都含某个所列基区间，其中有 \(|f'|=1\) 的点。因此 \(f'\) 在 \(C\) 的每一点都不连续。端点处使用单侧导数即可。

第（2）问取一个正测度疏朗完全集 \(C\) 即可。例如在 \([0,1]\) 的第 \(k\) 步，从每个保留区间的中央删去长 \(4^{-k}\) 的开段，共删去 \(\sum_{k\geq1}2^{k-1}4^{-k}=1/2\)。极限集无内点、无孤立点且测度为 \(1/2\)。代入上述构造，导函数间断点集即为 \(C\)。

\end{document}
